Since 2019 the Netherlands, home of ASML, the only maker of extreme ultraviolet lithography machines, has aligned its export controls with US demands while extracting little in return, as has Japan since 2023. Yet both resist some demands, such as limits on servicing their installed base in China. In our game-theoretic model, a security-dependent chokepoint holder bargains with its patron while the patron's rival bids for it. We prove each result and check it numerically. The holder's exit premium, the rival's credible offer net of the credible costs of exit, is a floor on what it obtains. Without rupture risk or agenda power, and without hold-up, its leverage is the positive part of that premium. The patron's demands are bounded by the holder's slack, which all its conditional threats share. Conditionality is the price of concessions, and a wasting slack gives the patron a preventive motive to demand control. The sign of Europe's premium turns on China's credible promises and, if Washington's threat to withhold frontier AI is credible, on that instrument, which on our orders of magnitude weighs more than how conditional the US guarantee is.
